\begin{table}[htbp]
\centering
\caption{Summary Statistics by Cohort Group}
\label{tab:sumstats}
\begin{adjustbox}{max width=\textwidth}
\begin{tabular}{lccc}
\toprule
 & Draft-Eligible & Older Control & Age Placebo \\
 & (Born 1915--1922) & (Born 1905--1914) & (Born 1895--1904) \\
\midrule
\multicolumn{4}{l}{\textit{Panel A: Occupational Outcomes}} \\[3pt]
Occupational Score, 1930 & 0.4 & 14.9 & 24.1 \\
 & (2.6) & (12.4) & (11.7) \\[3pt]
Occupational Score, 1940 & 16.4 & 24.1 & 25.6 \\
 & (11.5) & (10.9) & (11.4) \\[3pt]
Occupational Score, 1950 & 23.5 & 24.6 & 24.2 \\
 & (12.7) & (13.0) & (13.5) \\[3pt]
$\Delta$ OccScore (1940--1950) & 7.04 & 0.51 & -1.39 \\
\midrule
\multicolumn{4}{l}{\textit{Panel B: Demographics (1940)}} \\[3pt]
Education (years) & 7.7 & 7.2 & 6.1 \\
White (\%) & 95.3 & 95.8 & 95.9 \\
Married (\%) & 24.3 & 76.2 & 89.7 \\
Farm Resident (\%) & 26.1 & 21.8 & 21.9 \\
In Agriculture (\%) & 19.1 & 18.2 & 18.6 \\
\midrule
Observations & 4,358,625 & 4,773,710 & 4,249,305 \\
\bottomrule
\end{tabular}
\end{adjustbox}
\begin{tablenotes}
\small
\item \textit{Notes:} Standard deviations in parentheses. The sample consists of men linked across the 1930, 1940, and 1950 full-count U.S. Censuses via the IPUMS Machine Learning Panel (MLP) crosswalk. Occupational income scores are from the IPUMS OCCSCORE variable, which assigns 1950 median income values to 1950 occupation codes. Education is coded from the 1940 Census EDUC variable (categorical, converted to approximate years). Draft-eligible men were born 1915--1922 (ages 18--25 in 1940); older control men were born 1905--1914 (ages 26--35 in 1940); age placebo men were born 1895--1904 (ages 36--45 in 1940).
\end{tablenotes}
\end{table}
